\documentclass[a4paper,11pt]{article}
\usepackage[a4paper,margin=2cm]{geometry}
\usepackage[T1]{fontenc}
\usepackage[utf8]{inputenc}
\usepackage[ngerman,english]{babel}
\usepackage{lmodern}
\usepackage{microtype}
\usepackage{xcolor}
\usepackage{parskip}
\usepackage{enumitem}
\usepackage[most]{tcolorbox}
\usepackage{titlesec}
\usepackage[hidelinks]{hyperref}
\definecolor{HeaderBlueDark}{RGB}{70,110,170}
\definecolor{RowBlueLight}{RGB}{235,243,255}
\definecolor{RowBlueDark}{RGB}{220,233,250}
\definecolor{SoftGray}{RGB}{90,90,90}
\setlength{\parindent}{0pt}
\setlist[itemize]{leftmargin=1.4em,itemsep=0.35em,topsep=0.35em}
\linespread{1.06}
\titleformat{\section}[block]{\Large\bfseries\color{HeaderBlueDark}}{}{0pt}{}
\titlespacing{\section}{0pt}{1.3em}{0.8em}
\newtcolorbox{exbox}{enhanced,breakable,colback=RowBlueLight,colframe=RowBlueDark,boxrule=0.4pt,arc=1.5mm,left=1.2mm,right=1.2mm,top=1mm,bottom=1mm,title={Beispiele \textnormal{(Examples)}},fonttitle=\bfseries,colbacktitle=RowBlueDark,coltitle=black}
\NewDocumentEnvironment{grammarentry}{mm+b}{%
\begin{tcolorbox}[enhanced,breakable,colback=white,colframe=HeaderBlueDark,boxrule=0.6pt,arc=1.5mm,left=1.5mm,right=1.5mm,top=1mm,bottom=1mm,colbacktitle=HeaderBlueDark,coltitle=white,fonttitle=\bfseries,title={#1\hfill\normalfont\footnotesize #2}]
#3
\end{tcolorbox}
}{}
\newcommand{\GermanText}[1]{\textbf{Deutsch: }#1\par}
\newcommand{\EnglishText}[1]{{\small\color{SoftGray}\textbf{English: }#1\par}}
\newcommand{\ExampleItem}[2]{\item \textbf{#1}\\{\small\color{SoftGray}#2}}
\title{Die Relativpronomenform \glqq denen\grqq{}\\\large Relative Pronoun Form \glqq denen\grqq{}}
\date{}
\begin{document}
\maketitle
\tableofcontents
\newpage
\section{Grammatische Funktion (Grammatical function)}
\begin{grammarentry}{denen}{Dativ Plural}
\GermanText{\textbf{Denen} ist die Dativform des pluralischen Relativpronomens. Das Bezugswort kann maskulin, feminin oder neutral sein; im Plural ist das Genus an der Form nicht sichtbar. Nach einer Dativpräposition wie \textbf{aus} bleibt der Dativ erhalten.}
\EnglishText{\textbf{Denen} is the dative plural form of the relative pronoun. The antecedent may be masculine, feminine, or neuter; gender is not visible in the plural form. A dative preposition such as \textbf{aus} requires the dative.}
\end{grammarentry}
\section{Bezug zum Bezugswort (Relation to the antecedent)}
\begin{grammarentry}{Genus und Numerus}{Gender and number}
\GermanText{Genus und Numerus richten sich nach dem Bezugswort. Der Kasus dieser Relativpronomenform richtet sich nach ihrer Funktion im Relativsatz.}
\EnglishText{Gender and number come from the antecedent. The case of this relative-pronoun form is determined by its function inside the relative clause.}
\end{grammarentry}
\section{Beispiele (Examples)}
\begin{exbox}\begin{itemize}
\ExampleItem{Die Menschen, denen ich geholfen habe, waren dankbar.}{The people whom I helped were grateful.}
\ExampleItem{Die Gründe, aus denen ich den Computer benutze, sind verschieden.}{The reasons for which I use the computer are different.}
\end{itemize}\end{exbox}
\section{Satzstellung und Zeichensetzung (Word order and punctuation)}
\begin{grammarentry}{Relativsatz}{Relative clause}
\GermanText{Der Relativsatz wird durch ein Komma eingeleitet; das finite Verb steht am Ende.}
\EnglishText{A comma introduces the relative clause; its finite verb stands at the end.}
\end{grammarentry}
\section{Merksatz (Memory rule)}
\begin{grammarentry}{Die Form erkennen}{Recognizing the form}
\GermanText{Bestimme zuerst das Bezugswort und dann die Funktion des Relativpronomens im Relativsatz. Die Funktion bestimmt den Kasus.}
\EnglishText{First identify the antecedent, then determine the relative pronoun's function inside the relative clause. Its function determines the case.}
\end{grammarentry}
\end{document}
